% ====== 基础中文LaTeX文档模板 ======
\documentclass[UTF8]{article}
\usepackage{amsmath, amssymb, amsthm}  % 数学公式与定理环境
\usepackage{ctex}                      % 中文支持
\usepackage{geometry}                  % 页面布局调整
\usepackage{hyperref}                  % 超链接支持
\usepackage{graphicx}                  % 插入图片
\usepackage{listings}                  % 代码环境
\usepackage{xcolor}                    % 颜色定义
\usepackage{enumitem}                  % 列表格式定制

% ====== 基础设置 ======
\title{2026寒假数学建模期末考试答题卡}
\author{夏同~202530451676}




% ====== 文档正文 ======
\begin{document}
\maketitle

\subsection*{问题一}
\textbf{可用求解方法（至少三种）：}
\begin{enumerate}[label=(\arabic*), leftmargin=2em]
    \item \textbf{分支定界法}：系统搜索解空间，通过剪枝提高效率。
    \item \textbf{动态规划法}：用于有最优子结构的精确求解，如 Held-Karp 算法。
    \item \textbf{启发式/元启发式算法}：如遗传算法、模拟退火、蚁群算法，用于大规模问题近似求解。
\end{enumerate}

\textbf{图论讨论：}
\vspace{0.5em}

(1) \textbf{中国邮递员问题}：目标是寻找包含图中所有边的 \textbf{最短闭迹（回路）}。
\begin{itemize}[leftmargin=1.5em]
    \item \textbf{图模型}：将街道视为无向图 $G=(V,E)$ 的边，路口为顶点，边权表示长度或时间。
    \item \textbf{核心条件}：若 $G$ 为欧拉图（所有顶点度数为偶），则任一欧拉回路即为解；否则需通过添加重复边使所有顶点度数为偶，转化为欧拉图问题（“奇偶点图上作业法”）。
    \item \textbf{求解思路}：先确定所有奇度顶点，构造完全图并求其最小权完美匹配，将匹配边在原图中重复，使图变为欧拉图，再求欧拉回路。
\end{itemize}

(2) \textbf{旅行商问题}：目标是寻找访问图中所有顶点恰好一次的 \textbf{最短哈密顿回路}。
\begin{itemize}[leftmargin=1.5em]
    \item \textbf{图模型}：城市为顶点集 $V$，城市间距离为边权，通常建模为完全无向/有向加权图。
    \item \textbf{核心困难}：是 NP-hard 问题，顶点数 $n$ 的精确解搜索空间为 $(n-1)!/2$（对称情况）。
    \item \textbf{精确算法}：动态规划（Held-Karp，复杂度 $O(n^2 2^n)$）、分支定界。
    \item \textbf{近似算法}：最近邻法、最小生成树法（Christofides 算法，对度量 TSP 有 1.5 倍近似比）。
\end{itemize}


\subsection*{问题二}
\vspace{0.5em}

(1) \textbf{二维热传导模型}

\begin{itemize}[leftmargin=1.5em]
    \item \textbf{物理背景}：描述二维区域（如薄板）内部温度随时间的变化规律，基于 Fourier 热传导定律。
    \item \textbf{方程表达}：
    \begin{equation}
        \rho c \frac{\partial u}{\partial t} = \nabla \cdot (k \nabla u) + Q(x,y,t)
    \end{equation}
    其中 $\nabla = \left( \frac{\partial}{\partial x}, \frac{\partial}{\partial y} \right)$。若热导率 $k$ 为常数且无内热源（$Q=0$），则简化为：
    \begin{equation}
        \frac{\partial u}{\partial t} = \alpha \left( \frac{\partial^2 u}{\partial x^2} + \frac{\partial^2 u}{\partial y^2} \right), \quad \alpha = \frac{k}{\rho c}
    \end{equation}
    \item \textbf{参数含义}：
    \begin{itemize}
        \item $u(x,y,t)$：温度场（单位：K 或 °C）。
        \item $\rho$：材料密度（kg/m$^3$）。
        \item $c$：比热容（J/(kg·K)）。
        \item $k$：热导率（W/(m·K)）。
        \item $\alpha$：热扩散系数（m$^2$/s）。
        \item $Q(x,y,t)$：内热源强度（W/m$^3$）。
    \end{itemize}
    \item \textbf{初始条件}：
    \begin{equation}
        u(x,y,0) = f(x,y), \quad (x,y) \in \Omega
    \end{equation}
    $f(x,y)$ 为初始温度分布。
    \item \textbf{边界条件}（常见三类，定义在区域边界 $\partial \Omega$ 上）：
    \begin{itemize}
        \item \textbf{Dirichlet 条件}：$u(x,y,t) = g(x,y,t)$，给定边界温度。
        \item \textbf{Neumann 条件}：$k \frac{\partial u}{\partial \mathbf{n}} = q(x,y,t)$，给定边界热流密度（$\mathbf{n}$ 为外法向）。
        \item \textbf{Robin 条件}：$k \frac{\partial u}{\partial \mathbf{n}} + h(u - u_{\text{env}}) = 0$，对流换热（$h$ 为换热系数，$u_{\text{env}}$ 为环境温度）。
    \end{itemize}
\end{itemize}

(2) \textbf{二维弦振动模型}

\begin{itemize}[leftmargin=1.5em]
    \item \textbf{物理背景}：描述二维弹性薄膜（如鼓膜）在张力作用下的横向振动。
    \item \textbf{方程表达}：
    \begin{equation}
        \rho \frac{\partial^2 w}{\partial t^2} = T \nabla^2 w + f(x,y,t)
    \end{equation}
    其中 $\nabla^2 = \frac{\partial^2}{\partial x^2} + \frac{\partial^2}{\partial y^2}$ 为 Laplace 算子。若外力 $f=0$，则简化为二维波动方程：
    \begin{equation}
        \frac{\partial^2 w}{\partial t^2} = c^2 \left( \frac{\partial^2 w}{\partial x^2} + \frac{\partial^2 w}{\partial y^2} \right), \quad c = \sqrt{\frac{T}{\rho}}
    \end{equation}
    \item \textbf{参数含义}：
    \begin{itemize}
        \item $w(x,y,t)$：横向位移（单位：m）。
        \item $\rho$：面密度（kg/m$^2$）。
        \item $T$：张力（N/m）。
        \item $c$：波速（m/s）。
        \item $f(x,y,t)$：外力面密度（N/m$^2$）。
    \end{itemize}
    \item \textbf{初始条件}：需同时给出初始位移与初始速度：
    \begin{align}
        w(x,y,0) &= \phi(x,y), \\
        \frac{\partial w}{\partial t}(x,y,0) &= \psi(x,y), \quad (x,y) \in \Omega
    \end{align}
    \item \textbf{边界条件}（常见于固定边界或自由边界）：
    \begin{itemize}
        \item \textbf{Dirichlet 条件}：$w(x,y,t) = 0$（固定边界，位移为零）。
        \item \textbf{Neumann 条件}：$\frac{\partial w}{\partial \mathbf{n}} = 0$（自由边界，法向斜率为零）。
        \item \textbf{混合条件}：也可能出现弹性支撑等 Robin 类型条件。
    \end{itemize}
\end{itemize}




\subsection*{问题三}
\vspace{0.5em}

(1) \textbf{遗传算法 (Genetic Algorithm, GA) 步骤} \\
\textbf{核心思想}：模拟生物进化中的自然选择、交叉与变异机制。\\

\begin{enumerate}[label=\textbf{步骤 \arabic*：}, leftmargin=2em, align=left]
    \item \textbf{初始化种群} \\
    随机生成 $N$ 个个体构成初始种群 $P(0)$，每个个体（染色体）表示为解空间的一个编码（如二进制串、实数向量）。
    
    \item \textbf{适应度评价} \\
    计算每个个体 $x_i$ 的适应度 $f(x_i)$（目标函数值，需转换为非负且越大越好）。
    
    \item \textbf{选择（Selection）} \\
    根据适应度按概率选择父代个体。常用方法：
    \begin{itemize}
        \item 轮盘赌选择：选择概率 $p_i = f(x_i) / \sum_j f(x_j)$。
        \item 锦标赛选择：随机选取 $k$ 个个体，保留适应度最高者。
    \end{itemize}
    
    \item \textbf{交叉（Crossover）} \\
    以概率 $p_c$ 对选中的父代配对进行交叉操作，产生子代。例如：
    \begin{itemize}
        \item 单点交叉：随机选择一点，交换该点后的片段。
        \item 算术交叉：对实数编码，$ \text{child} = \alpha \cdot \text{parent}_1 + (1-\alpha) \cdot \text{parent}_2$。
    \end{itemize}
    
    \item \textbf{变异（Mutation）} \\
    以较小概率 $p_m$ 对子代基因进行随机变动（如二进制编码的位翻转），增加种群多样性。
    
    \item \textbf{更新种群} \\
    用新生成的子代替换当前种群（或采用精英保留策略保留最优个体），得到新一代种群 $P(t+1)$。
    
    \item \textbf{终止判断} \\
    若满足停止条件（如达到最大迭代次数、适应度收敛等），则输出最优解；否则返回步骤 2。
\end{enumerate}

(2) \textbf{模拟退火算法 (Simulated Annealing, SA) 步骤} \\
\textbf{核心思想}：模拟固体退火过程，通过控制温度参数，以一定概率接受劣解，避免陷入局部最优。\\

\begin{enumerate}[label=\textbf{步骤 \arabic*：}, leftmargin=2em, align=left]
    \item \textbf{初始化} \\
    设置初始温度 $T_0$（足够高），初始解 $x_0$，当前解 $x = x_0$，最优解 $x_{\text{best}} = x_0$，迭代计数器 $k=0$。
    
    \item \textbf{生成新解} \\
    在当前解 $x$ 的邻域中随机扰动产生新解 $x'$（如微小随机位移）。
    
    \item \textbf{计算能量差} \\
    计算目标函数值的增量 $\Delta E = f(x') - f(x)$（求最小化时，若 $\Delta E < 0$ 则新解更优）。
    
    \item \textbf{Metropolis 接受准则} \\
    决定是否接受新解 $x'$ 为当前解：
    \begin{itemize}
        \item 若 $\Delta E < 0$，则接受 $x'$（$x \leftarrow x'$）。
        \item 若 $\Delta E \geq 0$，则以概率 $p = \exp\left(-\dfrac{\Delta E}{T_k}\right)$ 接受 $x'$（否则保留 $x$）。
    \end{itemize}
    
    \item \textbf{更新最优解} \\
    若 $f(x') < f(x_{\text{best}})$，则更新 $x_{\text{best}} = x'$。
    
    \item \textbf{降温} \\
    按降温计划降低温度，常见方式：
    \begin{equation}
        T_{k+1} = \alpha \cdot T_k, \quad \alpha \in (0,1) \quad (\text{如 } \alpha=0.95)
    \end{equation}
    或 $T_{k+1} = \dfrac{T_0}{\ln(k+2)}$ 等。
    
    \item \textbf{终止判断} \\
    若满足停止条件（如温度低于阈值 $T_{\min}$、达到最大迭代次数、解长时间无改进），则输出 $x_{\text{best}}$；否则 $k \leftarrow k+1$，返回步骤 2。
\end{enumerate}


\subsection*{问题四}
\textbf{1. 点与点之间的距离}
\vspace{0.5em}

设 $x, y \in \mathbb{R}^n$，记 $x = (x_1, x_2, \dots, x_n)$, $y = (y_1, y_2, \dots, y_n)$。

\begin{itemize}[leftmargin=1.5em]
    \item \textbf{欧氏距离（$L_2$ 距离）}：最常用的几何距离
    \begin{equation}
        d(x, y) = \|x - y\|_2 = \sqrt{\sum_{i=1}^n (x_i - y_i)^2}
    \end{equation}
    
    \item \textbf{曼哈顿距离（$L_1$ 距离）}：网格路径距离
    \begin{equation}
        d_1(x, y) = \|x - y\|_1 = \sum_{i=1}^n |x_i - y_i|
    \end{equation}
    
    \item \textbf{切比雪夫距离（$L_\infty$ 距离）}：最大分量差
    \begin{equation}
        d_\infty(x, y) = \|x - y\|_\infty = \max_{1 \leq i \leq n} |x_i - y_i|
    \end{equation}
    
    \item \textbf{闵可夫斯基距离（$L_p$ 距离）}：统一形式
    \begin{equation}
        d_p(x, y) = \|x - y\|_p = \left( \sum_{i=1}^n |x_i - y_i|^p \right)^{1/p}, \quad p \geq 1
    \end{equation}
\end{itemize}

\textbf{案例}：
\begin{itemize}[leftmargin=1.5em]
    \item 在机器学习中，KNN 分类器常用欧氏距离度量样本相似度。
    \item 城市街区导航使用曼哈顿距离计算实际路径长度。
    \item 国际象棋中王的移动步数为切比雪夫距离。
\end{itemize}

\vspace{1em}
\textbf{2. 点与集合之间的距离}
\vspace{0.5em}

设点 $x \in \mathbb{R}^n$，集合 $A \subseteq \mathbb{R}^n$，且 $A \neq \emptyset$。

\begin{itemize}[leftmargin=1.5em]
    \item \textbf{点与集合的距离}定义为点与集合中所有点距离的下确界：
    \begin{equation}
        d(x, A) = \inf_{a \in A} d(x, a)
    \end{equation}
    若 $A$ 为闭集，下确界可改为最小值。
    
    \item 若采用欧氏距离，则：
    \begin{equation}
        d(x, A) = \inf_{a \in A} \|x - a\|_2
    \end{equation}
\end{itemize}

\textbf{案例}：
\begin{itemize}[leftmargin=1.5em]
    \item 在优化问题中，投影梯度法需要计算点到可行域的距离。
    \item 在模式识别中，判断一个样本点距离某类样本集的远近，用于异常检测。
    \item 计算几何中，求点到多边形的最短距离。
\end{itemize}

\vspace{1em}
\textbf{3. 集合与集合之间的距离}
\vspace{0.5em}

设 $A, B \subseteq \mathbb{R}^n$，且 $A, B \neq \emptyset$。常用定义有两种：

\begin{itemize}[leftmargin=1.5em]
    \item \textbf{豪斯多夫距离}：衡量两个集合间的最大不匹配程度
    \begin{align}
        d_H(A, B) &= \max \left\{ \sup_{a \in A} d(a, B), \sup_{b \in B} d(b, A) \right\} \\
        &= \max \left\{ \sup_{a \in A} \inf_{b \in B} d(a, b), \sup_{b \in B} \inf_{a \in A} d(a, b) \right\}
    \end{align}
    
    \item \textbf{最短距离}：两个集合中最近点对的距离
    \begin{equation}
        d_{\min}(A, B) = \inf_{\substack{a \in A \\ b \in B}} d(a, b)
    \end{equation}
    
    \item \textbf{质心距离}：两个集合质心间的距离
    \begin{equation}
        d_c(A, B) = \|\bar{a} - \bar{b}\|, \quad \bar{a} = \frac{1}{|A|} \sum_{a \in A} a, \quad \bar{b} = \frac{1}{|B|} \sum_{b \in B} b
    \end{equation}
\end{itemize}

\textbf{案例}：
\begin{itemize}[leftmargin=1.5em]
    \item 在图像处理中，豪斯多夫距离用于比较二值化图像的形状相似性。
    \item 在碰撞检测中，$d_{\min}(A, B) > 0$ 表示物体 $A$ 和 $B$ 未碰撞。
    \item 在聚类分析中，质心距离用于计算类间距离（如 K-means 聚类）。
    \item 在路径规划中，计算机器人（视为点集）与障碍物集的最小距离以保证安全。
\end{itemize}


\subsection*{问题五}
\textbf{1. 最大覆盖模型 (Maximum Coverage Problem, MCP)}
\vspace{0.5em}

\begin{itemize}[leftmargin=1.5em]
    \item \textbf{问题描述}：给定集合 $U = \{1, 2, \dots, m\}$（元素全集），$n$ 个子集 $S_1, S_2, \dots, S_n \subseteq U$，每个子集 $S_j$ 有覆盖权重（常为元素个数），以及预算 $k$（最多可选子集数量）。目标是选择至多 $k$ 个子集，使得被覆盖的元素总权重最大。
    
    \item \textbf{数学模型（0-1整数规划）}：
    \begin{align}
        \text{最大化} \quad & \sum_{i=1}^m w_i y_i \\
        \text{约束} \quad & \sum_{j=1}^n x_j \leq k \\
        & y_i \leq \sum_{j: i \in S_j} x_j, \quad \forall i = 1,\dots,m \\
        & x_j \in \{0,1\}, \quad y_i \in \{0,1\}, \quad \forall i,j
    \end{align}
    其中：
    \begin{itemize}
        \item $x_j = 1$ 表示选择子集 $S_j$，否则为 0。
        \item $y_i = 1$ 表示元素 $i$ 被至少一个选择的子集覆盖。
        \item $w_i$ 为元素 $i$ 的权重（常取 $w_i=1$）。
    \end{itemize}
    
    \item \textbf{求解特点}：该问题是 NP-hard，常用贪婪算法（每次选择覆盖新元素最多的子集）近似求解，近似比为 $1 - 1/e \approx 0.632$。
\end{itemize}

\textbf{案例}：
\begin{itemize}[leftmargin=1.5em]
    \item \textbf{基站选址}：$U$ 为需要覆盖的居民区，$S_j$ 为基站 $j$ 能覆盖的区域（有重叠），在建设预算限制下（最多建 $k$ 个基站），最大化覆盖的居民区数量。
    \item \textbf{广告投放}：$U$ 为潜在客户集合，$S_j$ 为平台 $j$ 能触达的客户群，在广告预算限制下选择 $k$ 个平台，最大化触达的客户数。
    \item \textbf{紧急设施选址}：在灾害救援中，选择 $k$ 个救援点，最大化能覆盖的受灾人口。
\end{itemize}

\vspace{1em}
\textbf{2. 集合覆盖模型 (Set Covering Problem, SCP)}
\vspace{0.5em}

\begin{itemize}[leftmargin=1.5em]
    \item \textbf{问题描述}：给定集合 $U = \{1, 2, \dots, m\}$，$n$ 个子集 $S_1, S_2, \dots, S_n \subseteq U$，每个子集 $S_j$ 有成本 $c_j > 0$。目标是选择一组子集，使得它们的并集等于 $U$（全覆盖），且总成本最小。
    
    \item \textbf{数学模型（0-1整数规划）}：
    \begin{align}
        \text{最小化} \quad & \sum_{j=1}^n c_j x_j \\
        \text{约束} \quad & \sum_{j: i \in S_j} x_j \geq 1, \quad \forall i = 1,\dots,m \\
        & x_j \in \{0,1\}, \quad \forall j = 1,\dots,n
    \end{align}
    其中 $x_j = 1$ 表示选择子集 $S_j$。
    
    \item \textbf{求解特点}：同样是 NP-hard，贪婪算法（每次选择性价比最高的子集，即“新覆盖元素数/成本”最大）的近似比为 $H(\max_j |S_j|) \leq 1 + \ln m$，其中 $H(d) = \sum_{i=1}^d 1/i$ 为调和数。
\end{itemize}

\textbf{案例}：
\begin{itemize}[leftmargin=1.5em]
    \item \textbf{航空公司机组调度}：$U$ 为所有需要执行的航班任务，$S_j$ 为一名飞行员能完成的一个合法班次（若干连续航班），成本为班次费用。选择最少成本的班次集合，覆盖所有航班。
    \item \textbf{电路板测试点选择}：$U$ 为所有待检测的故障类型，$S_j$ 为测试点 $j$ 能检测到的故障集合，成本为设置测试点的开销。选择最少数量的测试点覆盖所有故障类型。
    \item \textbf{公共服务设施选址}：$U$ 为所有社区，$S_j$ 为设施 $j$ 能服务的社区（如医院服务半径内），成本为建设费用。选择最小成本的设施集合，使每个社区至少被一个设施覆盖。
    \item \textbf{新闻摘要生成}：$U$ 为文档中所有关键信息点，$S_j$ 为句子 $j$ 包含的信息点集合。选择最少数量的句子，覆盖所有关键信息。
\end{itemize}

\vspace{0.5em}
\textbf{两种模型的主要区别}：
\begin{itemize}[leftmargin=1.5em]
    \item \textbf{目标}：最大覆盖模型在预算限制下最大化覆盖量；集合覆盖模型要求全覆盖下最小化成本。
    \item \textbf{约束}：最大覆盖有子集数量限制；集合覆盖有全覆盖要求。
    \item \textbf{应用场景}：最大覆盖用于资源有限时的部分覆盖优化；集合覆盖用于必须满足所有需求的最小成本方案。
\end{itemize}



\subsection*{问题六}

这是一个非常经典的数学建模问题（原型为1998年全国大学生数学建模竞赛A题）。为了给该公司设计出最合理的投资组合方案，我们需要通过建立数学模型，对“收益最大化”和“风险最小化”这两个相互制约的目标进行量化分析，寻找最佳平衡点（即帕累托有效前沿上的拐点）。以下是完整的建模、求解与方案设计过程：\\
一、 变量定义与数据预处理\\
1. 净收益率的计算\\
设 $x_i \ (i=1,2,3,4)$ 为投资于资产 $S_i$ 的资金数额；$x_0$ 为存入银行的资金数额。由于 $r_i$ 是平均收益率，$p_i$ 是交易费率，因此各项资产的实际净收益率 $R_i = r_i - p_i$：$S_1$: $R_1 = 28\% - 1\% = 27\%$$S_2$: $R_2 = 21\% - 2\% = 19\%$$S_3$: $R_3 = 23\% - 4.5\% = 18.5\%$$S_4$: $R_4 = 25\% - 6.5\% = 18.5\%$$S_0$ (存银行): $R_0 = 5\%$（无风险无交易费）\\
2. 交易费的合理简化\\
题干指出“当购买额不超过给定值 $u_i$ 时，交易费按购买计算（结合常规模型，可理解为实际投资规模远大于阈值 $u_i$ 时，交易费用直接呈比例 $p_i x_i$）”。观察数据可知，$u_i$ 最大值仅为 198 元，相对于总资金 10000 元极小。在实际操作中，只要决定投资某资产，投资额通常远大于 $u_i$（后面的计算结果也会印证这一点）。因此，为了使模型更加精炼，我们直接假设投资各资产的交易费完全按 $p_i x_i$ 比例扣除。\\
3. 总体风险度量题目指出总体风险用所投资的资产中最大的一个风险来度量。因此绝对风险值设为 $V = \max\{q_1 x_1, q_2 x_2, q_3 x_3, q_4 x_4\}$。\\
二、 数学模型的建立我们的目标是使得总净收益 $Q$ 最大，同时总体风险 $V$ 最小。建立如下双目标线性规划模型：
目标函数 1（最大化净收益 $Q$）：$$ \max Q = 0.27x_1 + 0.19x_2 + 0.185x_3 + 0.185x_4 + 0.05x_0 $$
目标函数 2（最小化总体风险 $V$）：$$ \min V = \max\{0.025x_1, 0.015x_2, 0.055x_3, 0.026x_4\} $$
约束条件（资金总额与非负限制）：$$ x_0 + x_1 + x_2 + x_3 + x_4 = 10000 $$$$ x_i \ge 0 \quad (i=0, 1, 2, 3, 4) $$
三、 模型的转化与求解针对多目标优化问题，通常采用固定风险上限法（$\epsilon$-约束法）。我们将风险 $V$ 视为投资者可容忍的一个给定上限，即要求每个风险资产满足 $q_i x_i \le V$。这实际上给出了各项风险资产的容量上限：$x_1 \le V / 0.025 = 40V$$x_2 \le V / 0.015 \approx 66.67V$$x_3 \le V / 0.055 \approx 18.18V$$x_4 \le V / 0.026 \approx 38.46V$四项风险资产的总容量限制为 $\approx 163.31V$。我们将 $x_0 = 10000 - (x_1+x_2+x_3+x_4)$ 代入收益函数 $Q$，整理得到边际收益方程：$Q = 500 + 0.22x_1 + 0.14x_2 + 0.135x_3 + 0.135x_4$，由于各项资产相对于银行存款的边际收益率均大于 0，最优策略是贪心算法：按照边际收益率从高到低（$S_1 > S_2 > S_4 = S_3$），在风险允许的范围内将资金尽量填满优势资产。根据风险限额 $V$ 的放宽，投资策略会出现三个关键拐点：

第一拐点（闲置资金清零点）：$V \approx 61.23$当 $V \le 61.23$ 时，即使把 $S_1 \sim S_4$ 全部买到风险上限，总容量依然不足 10000 元，剩余资金只能存银行。此时 $Q \approx 500 + 25.78V$。当 $V=61.23$ 时，四项风险资产恰好把 10000 元全部吸干（银行存款降为0）。此区间内，投资者每多承担 1 元绝对风险，总收益能暴增 25.78 元，性价比极高。

第二拐点（优势资产集中点）：$V = 93.75$当 $V > 61.23$ 以后，因为不再有闲置资金，想要进一步增加高收益资产（$S_1, S_2$）的投入，就必须从低收益资产（$S_3, S_4$）中抽调资金。当 $V = 93.75$ 时，$S_1$ 和 $S_2$ 的容量上限相加刚好达到 10000 元，此时 $S_3, S_4$ 被完全淘汰。在这个过渡区间，$Q \approx 1850 + 3.73V$。此区间内，每多承担 1 元风险，收益仅增加 3.73 元，性价比出现断崖式下跌。

第三拐点（极度激进梭哈点）：$V = 250$继续放宽风险，$S_2$ 也会被逐步抛弃，全资金涌向单项收益最高的 $S_1$。当 $V=250$ 时，10000元全额满仓 $S_1$。区间收益 $Q = 1900 + 3.2V$。💡 高级理论认知（驳斥“劣势资产无用论”）：表面上看，$S_4$ 净收益与 $S_3$ 相同，但风险却低一半，且它俩的收益/风险都逊于 $S_2$。许多常规思路会直接将 $S_3, S_4$ 作为“劣势资产”剔除。然而，由于单项资产容量受限，当决策者要求将风险控制在极低水平（例如 $V<93.75$）时，$S_1, S_2$ 根本吃不下 10000 元资金。此时为了不让资金烂在银行（仅5\%收益），依然需要将其投向 $S_3, S_4$（可赚18.5\%）。因此在稳健型投资中，$S_3, S_4$ 是不可或缺的有效拼图。\\

四、 最终投资组合方案推荐结合上述数学推理，“风险-收益边际性价比”在 $V=61.23$ 处出现了最显著的断层下跌。这也是数学意义上用最低限度风险撬动最高绝对资本使用率的黄金分割点。综合考量“净收益尽可能大，总体风险尽可能小”的题意，针对不同公司经营风格，推荐以下最优方案及备选：🌟 核心推荐方案：极致均衡稳健型（位于第一拐点，最强推荐）资金分配：购买资产 $S_1$： 2449.3 元购买资产 $S_2$： 4082.2 元购买资产 $S_3$： 1113.3 元购买资产 $S_4$： 2355.2 元存入银行 $S_0$： 0 元方案成效：预期净收益最大化： 2078.6 元 （综合年化回报率高达 20.79\%）。

总体风险最小化： 最大可能损失严控在 61.23 元 （占总资金风险率仅 0.61\%）。推荐理由： 这是一套极度分散化的高效抗风险方案。四种资产承担的绝对风险值完美均摊（均为 61.23），一分钱都不存银行，彻底榨干了所有优质资产在低风险阶段的“红利期”，是商业实操中最高性价比的组合。(且单项最少投资 1113.3 元也远大于门槛 $u_i$，完全符合现实可行性。)\\备选方案：精简进取型（位于第二拐点，供风险偏好较高的企业选择）资金分配：购买资产 $S_1$： 3750 元购买资产 $S_2$： 6250 元其他资产及存银行：均为 0 元方案成效：预期净收益： 2200 元 （综合回报率 22.00\%）。总体风险： 总体风险小幅升至 93.75 元 （风险率约 0.94\%）。推荐理由： 剔除了效率垫底的 $S_3, S_4$，资产管理更聚焦。相比核心方案，公司多承担了约 32.5 元的整体亏损风险，换取了额外 121 元的确定性收益增量，在合理的金融风险控制内（总体风险依然低于1%）。



\section*{第7题}

为了使比赛安排最为公平、公正、合理，我们需要把 14 个比赛项目排成一个时间序列，使得在相邻两个项目中“出现相同运动员”的总人次尽可能少。该问题可以转化为图论中的旅行商路径问题（Traveling Salesperson Path Problem, TSPP），即将各个项目视为“节点”，两个项目之间共同参赛的人数视为节点间的“边权”，目标是寻找一条经过所有节点且总权重最小的哈密顿路径。以下是完整的数学建模、求解与最佳赛程安排方案：\\
一、 符号说明与数据预处理\\
1. 参量与变量定义\\
比赛项目集合：$J = \{1, 2, \dots, 14\}$，共 14 个节点。运动员集合：$I = \{1, 2, \dots, 40\}$，共 40 名运动员。报名矩阵 $A$：定义一个 $40 \times 14$ 的 0-1 矩阵 $A = (a_{ij})$。如果第 $i$ 名运动员参加了第 $j$ 项比赛，则 $a_{ij} = 1$；否则 $a_{ij} = 0$。2. 冲突权重矩阵 $C$ 的构建为了衡量若某两项比赛“紧挨着举行”会造成的体力冲突情况，定义 $14 \times 14$ 的对称冲突矩阵 $C = (c_{jk})$。其中 $c_{jk}$ 表示同时报名了第 $j$ 项和第 $k$ 项比赛的运动员总人数。计算公式为：$$c_{jk} = \sum_{i=1}^{40} a_{ij} a_{ik} \quad (j, k \in J, j \neq k)$$规定主对角线元素 $c_{jj} = 0$。\\
二、 数学模型的建立\\
利用整数线性规划（ILP）思想建立赛程排序模型：\\
1. 决策变量引入 $0-1$ 决策变量 $x_{jk}$：$$ x_{jk} = \begin{cases} 1, & \text{在排表中，比赛 } k \text{ 被紧挨着安排在比赛 } j \text{ 之后举行} \\ 0, & \text{其他情况} \end{cases} $$
引入连续变量 $y_j \in [1, 14]$：表示比赛项目 $j$ 在整个赛事序列中的出场顺位。\\
2. 目标函数使连续参加两项比赛的运动员总人次（即路径的总权重） $Z$ 最小化：
$$ \min Z = \sum_{j=1}^{14} \sum_{k=1, k\neq j}^{14} c_{jk} \cdot x_{jk} $$
3. 约束条件赛程必须是一条不分叉、不闭合的单链时间线：出度与入度约束（每个项目最多只有一场直接前序和一场直接后序）：
$$ \sum_{k=1, k \neq j}^{14} x_{jk} \le 1 \quad (\forall j \in J) $$
$$ \sum_{j=1, j \neq k}^{14} x_{jk} \le 1 \quad (\forall k \in J) $$
路径长度约束（14 个项目首尾相连，正好包含 13 个相邻的时间间隔）：
$$ \sum_{j=1}^{14} \sum_{k=1, k \neq j}^{14} x_{jk} = 13 $$
MTZ消除子环约束（确保所有14个项目排成唯一连贯的赛程表，避免发生局部死循环排表）：$$ y_j - y_k + 14 \cdot x_{jk} \le 13 \quad (\forall j, k \in J, j \neq k) $$
三、 模型求解与最优排赛方案将上述矩阵与约束输入计算机求解器（或使用状态压缩动态规划求解）。令人振奋的是，该模型的最小目标函数值为 $Z_{min} = 0$。这意味着，通过科学排表，本届运动会完全可以做到没有任何一名运动员面临“连场作赛”的困境。其中一条完美符合要求的最优项目比赛出场顺序（由第1场至第14场）如下：项目 2 $\rightarrow$ 项目 5 $\rightarrow$ 项目 3 $\rightarrow$ 项目 1 $\rightarrow$ 项目 10 $\rightarrow$ 项目 9 $\rightarrow$ 项目 12 $\rightarrow$ 项目 6 $\rightarrow$ 项目 4 $\rightarrow$ 项目 13 $\rightarrow$ 项目 8 $\rightarrow$ 项目 7 $\rightarrow$ 项目 11 $\rightarrow$ 项目 14\\
四、 方案有效性验证
为了印证该赛程排表的绝对公平和科学性，我们随意抽查连排的场次进行手工验算：\\
第 1 场(项目2) 与 第 2 场(项目5)：报名项目 2 的运动员有 (1, 3, 9, 10, 11, 20, 25, 28, 29, 34)；报名项目 5 的运动员有 (6, 17, 26, 27, 30, 31, 35, 38, 39)。两组人员完全没有重叠，无一人需连赛。

第 8 场(项目6) 与 第 9 场(项目4)：项目 6 参赛者为 (6, 10, 22, 24, 32, 40)；项目 4 参赛者为 (3, 9, 13, 14, 30, 33, 34, 35, 36)。两组人员没有交集。

查验比赛任务最繁重的“24号运动员”：他报名了 4 个项目：6、7、13、14。在我们规划的时间表中，这四场比赛分别排在第8场、第12场、第10场、第14场。这名运动员每次出赛之间都能保证获得足够的休息场次。

结论：该数学模型准确寻找出了全局最优组合路径，达成了“连续参赛人数最小化（降为绝对最小值0）”的要求，实现了完全合理的比赛排程。

\section*{第8题}

\subsection*{1. 数据准备与预处理}

\begin{lstlisting}[language=Matlab,caption=数据准备与预处理]
% 载入数据
load fisheriris;
X = meas;
Y = species;

% 转换为数值与类别格式
X = double(X);
Y = categorical(Y);

% 划分数据集
train_idx = 1:100;
val_idx = 101:130;
test_idx = 131:150;

X_train = X(train_idx, :);
X_val = X(val_idx, :);
X_test = X(test_idx, :);
Y_train = Y(train_idx);
Y_val = Y(val_idx);
Y_test = Y(test_idx);

% 标准化处理
mu = mean(X_train);
sigma = std(X_train);
X_train = (X_train - mu) ./ sigma;
X_val = (X_val - mu) ./ sigma;
X_test = (X_test - mu) ./ sigma;
\end{lstlisting}

\subsection*{2. 网络结构设计}

\begin{lstlisting}[language=Matlab,caption=网络结构设计]
layers = [
    featureInputLayer(4, 'Name', 'input')
    fullyConnectedLayer(16, 'Name', 'fc1')
    reluLayer('Name', 'relu1')
    fullyConnectedLayer(8, 'Name', 'fc2')
    reluLayer('Name', 'relu2')
    fullyConnectedLayer(3, 'Name', 'fc3')
    softmaxLayer('Name', 'softmax')
    classificationLayer('Name', 'classoutput')
];
\end{lstlisting}

\subsection*{3. 训练配置与执行}

\begin{lstlisting}[language=Matlab,caption=训练配置与执行]
options = trainingOptions('adam', ...
    'MaxEpochs', 100, ...
    'MiniBatchSize', 16, ...
    'ValidationData', {X_val, Y_val}, ...
    'ValidationFrequency', 30, ...
    'Verbose', 1, ...
    'Plots', 'training-progress', ...
    'ValidationPatience', 5);

net = trainNetwork(X_train, Y_train, layers, options);
\end{lstlisting}

\subsection*{4. 模型评估}

\begin{lstlisting}[language=Matlab,caption=模型评估]
% 预测与准确率计算
YPred = classify(net, X_test);
accuracy = sum(YPred == Y_test) / numel(Y_test);
fprintf('测试集准确率: %.2f%%\n', accuracy * 100);

% 混淆矩阵
confMat = confusionmat(Y_test, YPred);
disp('混淆矩阵:');
disp(confMat);

% 预测结果表格
results = table(Y_test, YPred,'VariableNames',{'TrueLabel', 'PredLabel'});
disp('测试集预测结果:');
disp(results);

% 保存模型
save('iris_net.mat', 'net');
\end{lstlisting}

\subsection*{结果说明}
模型在训练过程中通过验证集监控性能，最终在独立测试集上评估分类准确率，并输出混淆矩阵与详细预测结果。模型权重保存至 \texttt{iris\_net.mat} 以便后续调用。

\section*{第9题}
\subsection*{综述}
首先确定游戏目标为结余更多资金且成功生存，在此前提下针对天气状况已知情形建立基于动态规划的多阶段决策模型；针对天气状况未知情形建立基于马尔可夫决策过程(MDP)的行动决策模型；在多个玩家共同进行游戏时，引入博弈的思想，进行良性与恶性竞争的博弈分析，建立基于互利共赢的行动决策模型。给出不同情形下对于一般玩家的最优策略。

问题一要求在只有一名玩家且已知未来30天天气的情况下，遵循游戏的规则，合理规划食物和水资源的购买与使用，给出能够得到最多收益的行动策略。玩家的行动以天为单位在时间上离散，可以分为每一个阶段。为达到游戏的最终目标需灵活规划每一个阶段的行动、路线、资源分配至最优，因此建立基于动态规划的多阶段决策模型，分析模型总结得到一般情况下玩家的最佳策略。使用算法对模型进行求解，得到第一二关的最优决策，将得到的最优决策与给出的一般性最佳策略进行对比验证，分析判断一般性最佳决策的可信度。

第二问要求在只有一名玩家且不知道未来的天气，只知道当前天气的条件下，遵循游戏的规则，合理规划食物和水资源的购买与使用，安排玩家每日的活动行程，给出能够得到最多收益的行动策略。与第一问相似，游戏目的亦为“得到最多的结余资金”，且游戏过程中的决策过程亦为离散的分阶段决策。由于游戏过程中的天气不可提前预知，仅知当前的天气信息，玩家需要根据当前已知的天气信息与身上剩余的物资数目，做出下一步行动的决策。因此可类比马尔可夫决策过程(MDP)，引入天气变量与动作值函数，结合问题一中的决策模型，建立基于马尔可夫决策过程的行动策略模型。分析模型总结得到一般情况下玩家的最佳策略，最后使用马尔可夫预测给出游戏期间内的天气，对模型进行仿真求解。最后将得到的策略与给出的一般性最佳策略进行对比验证，分析判断一般性最佳策略的可信度。

第三问要求在有多名玩家共同进行游戏时，分为：知晓未来时间的所有天气信息，且在游戏开始前制定好行动策略，后期的行动均按照制定好的方案进行。不提前制定好具体的行动方案只知晓当天的天气状况，且知晓当天其他玩家的行动方案与剩余物资的情况。这两种情况，分别给出对于一般玩家的行动准则。对于本题存在多个玩家同时进行游戏的情况，引入博弈的思想，将玩家间的竞争分为良性竞争与恶性竞争，发现恶性竞争会导致所有玩家亏损，而良性竞争会为所有玩家带来盈利。于是在两种已知条件不同的情况下，分别给出一般情况下玩家应采取的策略，再对第五第六关的实际情况进行分析，验证策略的可行性。
\subsection*{问题描述}
1.1 问题背景\\
考虑如下的小游戏：玩家凭借一张地图，利用初始资金购买一定数量的水和食物（包括食品和其他日常用品），从起点出发，在沙漠中行走。途中会遇到不同的天气，也可在矿山、村庄补充资金或资源，目标是在规定时间内到达终点，并保留尽可能多的资金。游戏的基本规则如下：\\
（1）以天为基本时间单位，游戏的开始时间为第0天，玩家位于起点。玩家必须在截止日期或之前到达终点，到达终点后该玩家的游戏结束。\\
（2）穿越沙漠需水和食物两种资源，它们的最小计量单位均为箱。每天玩家拥有的水和食物质量之和不能超过负重上限。若未到达终点而水或食物已耗尽，视为游戏失败。\\
（3）每天的天气为“晴朗”、“高温”、“沙暴”三种状况之一，沙漠中所有区域的天气相同。\\
（4）每天玩家可从地图中的某个区域到达与之相邻的另一个区域，也可在原地停留。沙暴日必须在原地停留。\\
（5）玩家在原地停留一天消耗的资源数量称为基础消耗量，行走一天消耗的资源数量为基础消耗量的2倍。\\
（6）玩家第0天可在起点处用初始资金以基准价格购买水和食物。玩家可在起点停留或回到起点，但不能多次在起点购买资源。玩家到达终点后可退回剩余的水和食物，每箱退回价格为基准价格的一半。\\
（7）玩家在矿山停留时，可通过挖矿获得资金，挖矿一天获得的资金量称为基础收益。如果挖矿，消耗的资源数量为基础消耗量的3倍；如果不挖矿，消耗的资源数量为基础消耗量。到达矿山当天不能挖矿。沙暴日也可挖矿。\\
（8）玩家经过或在村庄停留时可用剩余的初始资金或挖矿获得的资金随时购买水和食物，每箱价格为基准价格的2倍。\\
1.2 需要求解的问题\\
对题目进行分析总结，得到各个问需要解决的问题：\\
问题一：假设只有一名玩家，在整个游戏时段内每天天气状况事先全部已知，试给出一般情况下玩家的最优策略。求解附件中的“第一关”和“第二关”，并将相应结果分别填入Result.xlsx。\\
问题二：假设只有一名玩家，玩家仅知道当天的天气状况，可据此决定当天的行动方案，试给出一般情况下玩家的最佳策略，并对附件中的“第三关”和“第四关”进行具体讨论。\\
问题三：现有n名玩家，有相同的初始资金，且同时从起点出发。若有k名玩家：\\
a. 通过相同的转移路径，且起始点相同，资源消耗量为基础的2k倍。\\
b. 同时在同一矿山挖矿，收益为基础收益的 $\frac1k$倍。\\
c. 同时在同一村庄购买资源，价格为原价格的4倍。\\
求解以下问题：\\
假设在整个游戏时段内每天天气状况事先全部已知，每名玩家的行动方案需在第三天确定且此后不能更改。试给出一般情况下玩家应采取的策略，并对附件中的“第五关”进行具体讨论。\\
假设所有玩家仅知道当天的天气状况，从第五天起，每名玩家在当天行动结束后均知道其余玩家当天的行动方案和剩余的资源数量，随后确定各自第二天的行动方案。试给出一般情况下玩家应采取的策略，并对附件中的“第六关”进行具体讨论。\\
\subsection*{规则量化}
每天的天气为晴朗、高温、沙暴三种状况之一，设：
\[
W_t = \begin{cases}
1 & \text{晴朗} \\
2 & \text{高温} \\
3 & \text{沙暴}
\end{cases}
\]
每天玩家可从地图中的某个区域到达与之相邻的另一个区域，也可在原地停留，沙暴日必须在原地停留。定义邻接矩阵：
\[
S_{ij} = \begin{cases}
1 & \text{区域 } i \text{ 与区域 } j \text{ 相邻} \\
0 & \text{不相邻}
\end{cases} \quad (4)
\]
前进标志：
\[
G_t = \begin{cases}
1 & \text{天气不为沙暴} \\
0 & \text{天气为沙暴}
\end{cases} \quad (5)
\]
则从第0天到第 \(t\) 天行走的路线 \(d_t\) 的迭代：
\[
d_{t+1} = d_t + G_t \cdot S_{ij} \quad (6)
\]
玩家在原地停留一天消耗的资源数量称为基础消耗量，行走一天消耗的资源数量为基础消耗量的2倍，则第 \(t+1\) 天时水和食物的剩余箱数为：
\[
\begin{cases}
n_{w(t+1)} = n_{wt} - (G_t + 1) \cdot BM_{wt} \\
n_{f(t+1)} = n_{ft} - (G_t + 1) \cdot BM_{ft}
\end{cases} \quad (7)
\]
玩家第0天可在起点处用初始资金以基准价格购买水和食物。玩家可在起点停留或回到起点，但不能多次在起点购买资源。玩家到达终点后可退回剩余的水和食物，每箱退回价格为基准价格的一半，即：
\[
\begin{array}{r}
n_{w0} \cdot P_w + n_{f0} \cdot P_f = M_0 - M_1 \\
M_T = M_{T-1} + n_{wt} \cdot \frac{P_w}{2} + n_{ft} \cdot \frac{P_f}{2}
\end{array} \quad (8)
\]
其中 \(n_{w0}, n_{f0}\) 为第0天购买的物资，\(n_{wt}, n_{ft}\) 为第 \(T\) 天剩余的物资。玩家在矿山停留时，可通过挖矿获得资金，挖矿一天获得的资金量称为基础收益。如果挖矿，消耗的资源数量为基础消耗量的3倍；如果不挖矿，消耗的资源数量为基础消耗量。到达矿山当天不能挖矿，沙暴日也可挖矿。玩家第 \(t\) 天所处的位置为：
\[
d_{t+1} - d_t = S_{ij} \quad (9)
\]
若 \(j = node_{Mp}\)（矿山），则：
\[
\begin{cases}
n_{w(t+1)} = n_{wt} - (2D_t + 1) \cdot BM_{wt} \\
n_{f(t+1)} = n_{ft} - (2D_t + 1) \cdot BM_{ft}
\end{cases} \quad (10)
\]
其中 \(D_t\) 为挖矿标志，挖矿时置1，否则0。玩家经过或在村庄停留时可用剩余的初始资金或挖矿获得的资金随时购买水和食物，每箱价格为基准价格的2倍。若 \(j = node_{Cp}\)（村庄），购买物资：
\[
\begin{cases}
n_{w(t+1)} = B_{wp} + n_{wt} \\
n_{f(t+1)} = B_{fp} + n_{ft}
\end{cases} \quad (12)
\]
第 \(t\) 天购买完物资后，第 \(t+1\) 天时现金剩余：
\[
M_{t+1} = M_t - 2B_{wp} \cdot P_w - 2B_{fp} \cdot P_f \quad (13)
\]
\subsection*{基于动态规划的行动决策模型}
由于玩家的行动以天为单位，在时间上离散地做出行动决策，每一步行动的结果都会对下一步的行动产生影响，同样每一步的行动结果都会对最终的收益产生影响。因此将该问题类比为离散时间下的确定型动态规划问题。由此建立基于离散时间下动态规划的行动决策模型，对玩家的行动策略进行求解。

阶段：指的是问题需要做出决策的步数。阶段数目常记为 \(n\)，对应着 \(n\) 个阶段的决策问题。阶段的序号记为 \(k\)，\(k = 1,2,\dots\)，称为阶段变量。

状态：各阶段开始时的客观条件称之为状态，第 \(k\) 阶段的状态常用状态变量 \(S_k\) 表示，状态变量的取值的集合称为状态集合。第1阶段的状态变量 \(S_1 = \{A\}\)，第2阶段的状态变量为 \(S_2 = \{B_1,B_2\}\)。

决策：从某个状态出发，在若干各不同的方案中做出的选择称之为决策。用于表示决策的变量 \(u_k(S_k)\) 称之为决策变量。\(u_k(S_k)\) 表示第 \(k\) 阶段当处于 \(S_k\) 状态时的决策变量。如图中的多阶段决策过程中，\(u_3(C_2) = D_1\) 表示走到 \(C\) 阶段，当处于 \(C_2\) 路口时，下一步的前加方向为 \(D_1\)。决策变量允许的取值范围称为允许决策集合，第 \(k\) 阶段状态为 \(S_k\) 时允许决策集合记为 \(D_k(S_k)\)，如图中 \(B_1\) 状态允许的决策集合为 \(D_2(B_1) = \{C_1,C_2,C_3\}\)。

指标函数：阶段指标函数是指第 \(k\) 阶段从状态 \(S_k\) 出发，采取决策 \(U_k\) 时的效益，使用
\[
\begin{cases}
f_k(S_k) = \underset{u_k\in D_k(S_k)}{\operatorname{opt}}\{V_k(S_k,u_k(S_k)) + f_{k+1}(u_k(S_k))\} \quad k = n,n-1,\dots,2,1 \\
f_{n+1}(S_{n+1}) = 0
\end{cases} \quad (14)
\]
其中 \(V_k(S_k,u_k)\) 为本次决策所产生的效益，opt可取max，min。

效益最大化目标：游戏规则要求，需要在游戏结束时，得到最多的收益，即身上的资金数目最高。首先考虑从出发点直接前往终点的路径，发现玩家身上的资金数额变化为负值，即没有得到收益，因此可以直接排除从起点出发直接前往终点的情况。观察游戏规则易知，整个游戏过程中唯一的盈利点为在矿山挖矿，因此玩家在游戏的开始应尽量以最快速度前往矿山，降低在路上的物资损耗，进行挖矿操作，赚取利润。

将第 \(k\) 阶段从状态 \(S_k\) 出发，采取决策 \(u_k\) 时的效益 \(f_k(S_{k+1})\) 表示为
\[
\begin{cases}
f_k(S_k) = \underset{u_k\in D_k(S_k)}{\operatorname{opt}}\{V_k(S_k,u_k(S_k)) + f_{k+1}(u_k(S_k))\} \quad k = n,n-1,\dots,2,1 \\
f_{n+1}(S_{n+1}) = 0
\end{cases} \quad (15)
\]
则该模型的最终目标即为
\[
f_k(S_k) = \underset{u_k\in D_k(S_k)}{\operatorname{max}}\{V_k(S_k,u_k(S_k)) + f_{k+1}(u_k(S_k))\} \quad k = n,n-1,\dots,2,1 \quad (16)
\]
其中，第 \(n+1\) 阶段的收益为0，即
\[
f_{n+1}(S_{n+1}) = 0 \quad (17)
\]

决策变量的确定：当物资大于矿山到村庄所需要的物资时，说明玩家可以继续在矿山进行挖矿的盈利活动，或者移动前往矿山进行挖矿，即
\[
u_k(C) = C, \quad u_k(A) = C, \quad u_k(B) = C \quad (18)
\]
当下一步操作后，剩余物资少于矿山到村庄所需要的物资数目时，玩家离开矿山，前往村庄进行补给，即
\[
u_k(C) = B \quad (19)
\]
假设从村庄回到终点所需的时间为 \(x_1\) 天，则当前时刻 \(t = T - x_1\) 时，从村庄前往终点，即
\[
u_{T-x_1+1}(B) = D \quad (20)
\]
假设从矿山回到终点所需的时间为 \(x_2\) 天，则当前时刻 \(t = T - x_2\)，离开矿山前往终点，即
\[
u_{T-x_2+1}(C) = D \quad (21)
\]

模型综合：结合以上论述，得到基于动态规划的决策模型如下：
\[
\begin{aligned}
f_k(S_k) &= \max_{u_k\in D_k(S_k)}\{V_k(S_k,u_k(S_k)) + f_{k+1}(u_k(S_k))\} \quad k = n,n-1,\dots,2,1 \\
f_{n+1}(S_{n+1}) &= 0 \\
u_k(C) &= \begin{cases}
C & \text{(进行下一步操作后剩余物资大于矿山到村庄所需物资)} \\
B & \text{(进行下一步操作后剩余物资小于从矿山到村庄所需物资)} \\
D & (t = T - x_2)
\end{cases} \\
u_k(B) &= \begin{cases}
D & (t = T - x_1) \\
C & \text{(进行下一步操作后剩余物资大于矿山到村庄所需物资)}
\end{cases} \\
u_k(A) &= C \quad \text{(进行下一步操作后剩余物资大于矿山到村庄所需物资)}
\end{aligned}
\]
通过分析上述模型建立过程中的内容，给出如下几条一般情况下玩家的策略：
\begin{enumerate}
    \item 购买满足全过程的食物，穿越沙漠过程中不在村庄购买食物。且剩余负重全部用来购买水。直至总负重达到允许的最大值，若购买的水不足以支撑从起点到村庄过程中的消耗则购买足量的水，其余空间购买食物。
    \item 在少数情况下，可根据实际情况，少带一些食物多带一些水来推迟补给的时间，延长挖矿的时间，并在后续的补给中补齐所需食物。
    \item 从起点到矿山路程中，不刻意停止，除非遇到沙暴天气。
    \item 在矿山挖矿期间，可在沙暴消耗量较高的天气里刻意休整一到两天，使利益最大化。
    \item 当挖矿期间所剩物资，仅支持从矿山前往村庄，则立即前往村庄补给物资。
    \item 二次补给结束后，前往终点方向的挖矿点或终点，根据具体地图和剩余时间选择进行二次挖矿或前往终点，并根据选择补给后续足够数目的物资。
\end{enumerate}

\subsection*{问题二：基于马尔可夫决策过程的行动决策模型}
对于本题来说，玩家的行为在时间上是离散的，其行为决策的做出只与当天的天气状况、身上剩余的物资数目和做出决策后的预期收益有关，即与只当前已知的信息有关，与过去的历史信息无关，因此可以认为本题的决策过程也具有马尔可夫性。

一个马尔可夫决策过程由一个四元组构成 \(M = (S, A, P_{sa}, R)\)，其中：
\begin{itemize}
    \item \(S\)：表示状态集(states)，有 \(s \in S\)，用 \(s_i\) 表示第 \(i\) 步的状态，即玩家进行第 \(i\) 天游戏时，玩家所处的位置状况。
    \item \(A\)：表示一组动作(actions)，有 \(a \in A\)，用 \(a_i\) 表示第 \(i\) 步的动作，即玩家进行第 \(i\) 天游戏时决定挖矿、行走、停留、买物资等行为。
    \item \(P_{sa}\)：表示状态转移的概率，用于表示在状态 \(s\) 下，经过动作 \(a\)，转移到其他状态的概率分布。如玩家在 \(s\) 的位置状态下进行动作 \(a\)，转移到 \(s'\) 位置状态的概率为 \(P(s'|s,a)\)。
    \item \(R\)：表示奖励函数，\(r(s'|s,a)\) 表示在 \(s\) 状态下，进行动作 \(a\) 后，状态转移为 \(s'\)，得到的奖励，即表示玩家在 \(s\) 的位置状态下，经过动作 \(a\) 后转移到位置状态 \(s'\)，并由此决策而产生的收益。
    \item \(V\)：表示值函数。对于之前的任意状态 \(s\) 和动作 \(a\)，立即奖励函数 \(R(s,a)\) 无法说明策略的好坏，因此定义值函数 \(V^\pi(s)\) 来表明当前状态下策略 \(\pi\) 的长期影响。\(V^\pi(s)\) 表示在策略 \(\pi = S \rightarrow A\) 下，状态 \(s\) 的值函数，表示当前状态下策略 \(\pi\) 的长期影响。使用如下公式：
    \[
    V^\pi(s) = E_\pi\left[\sum_{i=0}^{\infty} \gamma^i R_i \big|_{s_0 = s}\right] \quad (24)
    \]
    其中 \(\gamma \in [0,1]\) 为折合因子，表示了未来的回报对于当前回报的重要程度。当 \(\gamma\) 为0时，只考虑立即回报，不考虑短期回报；当 \(\gamma\) 为1时，长期回报与短期回报被同等看待。
\end{itemize}
奖励函数：
当玩家不在矿山时，奖励函数主要由在当天的天气下进行行动所消耗的物资数目，即
\[
R(\{S|S\neq S_{in}\}, M, m_w, m_f) = \sum_{i=1}^{2}(P_i \cdot \Delta M_i + C)
\]
其中 \(P_i\) 为第 \(i\) 种天气出现的概率，\(\Delta M_i\) 为第 \(i\) 种天气对应状态的转移资金消耗，\(C\) 为使得 \(\Delta M_i + C\) 为正数的常数。当进行某种行动消耗的物资越多时，奖励函数 \(R\) 的数值越小，则玩家会倾向于选择奖励函数更大的行动，不选择奖励函数较小的行动，使得在该状态下的玩家做出最优的行动决策。

当玩家位置在矿山时，当水或食物消耗量达到临界阈值时，将奖励函数置为0，迫使玩家做出离开矿山的行动决策，即
\[
R(S_{in}, M, m_w, m_f) = \sum_{i=1}^{2}(P_i \cdot \Delta M_i + C + BI \cdot D_i) \cdot \epsilon(m_w - m_w') \cdot \epsilon(m_f - m_f') 
\]
其中 \(m_w', m_f'\) 分别为水与食物质量的阈值，即为从该点前往村庄补给物资或前往终点所需要的消耗的物资；\(BI\) 为挖矿时的基础收益；\(D_i\) 为第 \(i\) 天时的挖矿标志，挖矿为1，不挖矿为0。

值函数：
值函数的定义如式(24)所示
\[
V^\pi(s) = E_\pi\left[\sum_{i=0}^{\infty} \gamma^i R_i \big|_{s_0 = s}\right] 
\]
在 \(V^\pi(s,a)\) 当中，策略 \(\pi\)，初始状态 \(s\) 均初始给定，而初始动作 \(a\) 是由策略 \(\pi\) 和状态 \(s\) 决定，即
\[
a = \pi(s) 
\]
给定策略 \(\pi\) 和初始状态 \(s\)，则动作 \(a = \pi(s)\)，下一个时刻将以概率 \(P(s'|s,a)\) 转向下一个状态 \(s'\)，那么 \(V^\pi(s)\) 可展开为：
\[
V^\pi(s) = \sum_{s' \in S} P(s'|s,a)[R(s'|s,a) + \gamma V^\pi(s')] 
\]
给定当前状态 \(s\) 和当前动作 \(a\)，在未来遵循策略 \(\pi\)，那么系统将以概率 \(P(s'|s,a)\) 转向下个状态 \(s'\)，我们可定义动作值函数 \(Q^\pi(s,a)\) 为
\[
Q^\pi(s,a) = \sum_{s' \in S} P(s'|s,a)[R(s'|s,a) + \gamma V^\pi(s')] 
\]
在 \(Q^\pi(s,a)\) 当中，策略 \(\pi\)，初始状态 \(s\)，当前的动作 \(a\) 均初始给定。

则MDP的最优策略可以由下式表示
\[
\pi^* = \arg \max V^\pi(s), \quad (\forall s) 
\]
即寻找的是在任意初始条件 \(s\) 下，能够最大化值函数的策略 \(\pi^*\)。

通过分析关卡四的模型参数可知，在矿山挖一天矿的收益为1000元，足够弥补在前往矿山的过程中所消耗的物资，因此如一般性准则，在动作值函数Q的指引下，玩家会先前往矿山。当玩家的物资第一次不足以支持继续挖矿时，时间依然充裕，根据7一般性准则，玩家会前往村庄进行补给，并返回矿山继续挖矿。当玩家的剩余游戏天数即将达到前往终点所需的时间阈值时，根据7.3中的一般性准则，玩家会停止挖矿，离开矿山，并在动作值函数Q的引导下前往终点。由此可以认为，一般性行动策略可信度较高。

\subsection*{基于博弈思想的多玩家决策竞争模型}
各玩家之间均存在着竞争，分为良性竞争与恶性竞争。

恶性竞争：
题目规则规定：
\begin{itemize}
    \item 当k名玩家通过相同的路径进行转移时，消耗量为基础消耗量的2k倍；
    \item 当k名玩家同时在矿山挖矿时，收益为基础收益的 \(\frac{1}{k}\)；
    \item 当k名玩家在同一村庄购买资源时，价格为基础价格的4倍。
\end{itemize}
在游戏过程中会存在以下恶性竞争：
\begin{enumerate}
    \item 每名玩家都为了赚取更多的资金，无计划地前往矿山进行挖矿，当天气恶劣时，由于基础收益的降低，挖矿获得的收益将不足以弥补物资的消耗量，此种竞争将导致所有玩家的收益降低，甚至亏损。
    \item 每名玩家为了更快地到达矿山进行挖矿、前往村庄进行补给、前往终点，无计划地走最短的转移路径，可能出现在同一时刻通过相同的路径进行位置的转移，将增加所有玩家的物资消耗。
    \item 每名玩家为了更快的前往村庄进行购物，可能出现玩家同时购买物资的情况，极大地增加玩家购买物资的资金消耗。
    \item 某名玩家为了迫使其他玩家尽快退出游戏而让自己得到独自在矿山挖矿的机会，在其他玩家前往商店购买物资时，自己只购买极少量的物资，来让其他玩家采购物资的花费极大地增加，而不得不提前退出游戏。该竞争方法会带来非常可观的收益，但其他玩家也有可能会想到该竞争方式，可能会使用该方式与自己进行竞争，因此存在极高的风险。
\end{enumerate}
以上几种竞争方式均会使得自己与其他玩家的收益同步降低，减少游戏结束时的剩余资金，因此可以认为是恶性竞争。由于在本题中，正面的直接竞争都将导致所有玩家的收益降低，与游戏的目的“得到更多的剩余资金”相悖，均属于恶性竞争。于是引入另外一种玩家间相互配合，互利共赢的良性竞争策略。

良性竞争：
\begin{enumerate}
    \item 对于游戏中获取资金最关键的挖矿行为，应尽量避免同时有多个人在矿山内挖矿，只允许能够使挖矿者获得正向收益的挖矿人数同时在矿山内挖矿，其余玩家在外等候。
    \item 对于游戏中前往村庄进行物资补充的过程，各名玩家均应避免与其他玩家在同一时刻购买物资。
    \item 对于游戏中位置迁移的过程，各玩家均应避免在同一时刻经过同一路径进行位置的迁移，可选择错时前行或选择不同的迁移路径。
    \item 在游戏过程中所有玩家均不使用恶性手段，使其他玩家蒙受巨额损失。
\end{enumerate}
对于玩家A来说，如果为了最大化自己的收益而选择恶性竞争，那么其他玩家也有可能选择恶性竞争来最大化自己的收益，则玩家A会有较大概率因为其他玩家发起的恶性竞争而降低自己的收益，甚至亏损。因此对于所有玩家来说，最优的策略是与其他玩家达成共识，只进行良性竞争，在使得每位玩家收益相近的条件下，追求自己的最大收益。则良性竞争的目标函数如下：

个人利益最大化目标，即最终持有的资金要更多：
\[
\max \sum_{i=1}^{n} M_{Ti} \quad (38)
\]
其中 \(M_{Ti}\) 为第 \(i\) 名玩家最终持有的资金数目。

各名玩家收益相近目标，即各名玩家最终持有资金数目方差要更小：
\[
\min \frac{\sum_{i=1}^{n}(M_j - \frac{\sum_{i=1}^{n} M_{Ti}}{n})^2}{n-1} 
\]
于是得到良性竞争的目标函数如下：
\[
\begin{cases}
\max \sum_{i=1}^{n} M_{Ti} \\
\min \frac{\sum_{i=1}^{n}(M_j - \frac{\sum_{i=1}^{n} M_{Ti}}{n})^2}{n-1}
\end{cases}
\]
一般性策略：
\begin{enumerate}
    \item 根据已知的天气信息，使用第一问的模型，分别计算每个玩家独立行动时的最优策略和收益。
    \item 当多名玩家在矿山挖矿时，根据天气情况，计算折损后的真实收益，若挖矿收益为正，则可安排多人同时挖矿，否则应错开时间。
    \item 为前往村庄购买物资留出余量，避免多人同时购买物资。
    \item 当需要从同一起点前往同一点时，多个玩家选择不同的转移路径，防止路上转移动作的消耗翻倍。当走不同的路径最终结果不优时，也可以选择走相同的转移路径。
    \item 多个玩家统一协调安排，使得各个玩家收益接近。
\end{enumerate}
情况二为天气未知情况下，策略未完全确定，在多名玩家共同进行游戏时，玩家按照当天天气情况，剩余物资情况，其他玩家的所处状态进行分析。

一般性策略：
总体规则：使用动作价值函数 \(Q\) 选择下一步的行动与目的地，其中
\[
Q^\pi(s,a) = \sum_{s' \in S} P(s'|s,a)[R(s'|s,a) + \gamma V^\pi(s')] 
\]
细化规则：
\begin{enumerate}
    \item 根据计算得到的各种天气出现的概率，预估未来的天气状况，结合所剩余的游戏天数和当前位置与终点的距离状况决定是否继续挖矿或前往终点。
    \item 根据身上剩余的物资数目，结合当前时刻的天气状况决定是否前往村庄进行补给；在补给后综合考虑剩余天数，预估未来的天气状况决定是否前往矿山进行挖矿或直接前往终点。
    \item 综合考虑矿山挖矿收益状况与前往矿山与在矿山挖矿所消耗的物资数目，判断是否盈利，而选择从起点出发后是否前往矿山，或者直接前往终点。
    \item 当多名玩家在矿山挖矿时，结合对未来天气的预测，计算折损后的真实收益，用以进行后续的判断。
    \item 为前往村庄购买物资留出余时间，避免多人同时购买物资。
    \item 当需要从同一起点前往同一条点时，多个玩家选择不同的转移路径，防止路上转移动作的消耗翻倍。当走不同的路径最终结果不优时，也可以选择走相同的转移路径。
    \item 多个玩家统一协调安排，使得各个玩家收益接近。
\end{enumerate}

模型运算结果和matlab代码在上传的另外一个压缩包中。
% --- 文档结束 ---
\end{document}